\chapter{Introduction\label{cha:chapter1}}

Large Language Models (LLMs) are widely used as assistants for software engineering work: planning, code generation, debugging, test generation/automation, documentation, and agentic tool workflows \cite{brown2020fewshot, ouyang2022instructgpt, chen2021codex, jimenez2023swebench, yao2023react}. In those systems, models' responses are passed to a parser, executed by a test harness, extracted into a function body, checked by an automated judge, or used as input to another tool. Therefore, not only correctness but also the surface behavior matters. A response may contain the right algorithm but still fail because the harness cannot parse it: it may include extra markdown, explanations, invalid JSON/XML, no clear final answer, or text that the evaluation pipeline cannot easily extract.

\section{Motivation}

Those settings make prompt framing an important part of system design. Developers often use techniques such as polite language, assigning expert roles, pressuring for accuracy, or formatting the prompt in JSON or XML \cite{li2023emotionprompt, yin2024politeness, liu2024structuredoutput}. Such common advice may be useful in some contexts, but evidence of the effectiveness of many prompt rules is not strong. Prior work shows that prompt wording, formatting, and ordering can change model behavior in different directions \cite{zhao2021calibrate, lu2021fantastically, sclar2023spurious, zhuo2024prosa}. A prompt that improves one model on one task may fail on another model or task. A prompt that helps a cloud model return parseable code may harm a local math model by encouraging compact answers. A role prompt may clarify the domain expectation for one model, but may make another model more verbose.

This led to two problems. The first problem is the reproducibility problem for empirical LLM research. If prompt phrasing changes model behavior in ways that are not reported or controlled, the resulting performance may be caused by prompt design rather than model ability \cite{hua2025promptartifact}. The second problem is the reliability problem for coding agents. If prompt choices are selected from folklore rather than experiment and validation, an engineering workflow may become slower, more expensive, less parsable, or produce more errors.

\section{Objective}

This thesis studies prompt framing sensitivity: how semantically similar prompts lead to different behaviors due to framing choices. We examine three factors: tone (neutral, polite, aggressive), role (no role and expert role), and input format (plaintext, JSON, XML). These three factors create an 18-category prompt matrix, and each base task is evaluated under every category, across coding and mathematical reasoning tasks.

\section{Thesis Scope}

The main software-engineering task investigated in this thesis is Python code generation. The model is asked to generate Python code satisfying given requirements; an evaluator then extracts the code from the response and evaluates it against a list of tests. The secondary domain is mathematical reasoning, where only the final answer will be extracted and evaluated. The distinction matters: code workflows often benefit from clean, parseable output, whereas in math only the final number matters, even though deriving it may require enough space for intermediate reasoning.

The primary result of this research is intentionally conservative. While the capacity of the model overshadows task performance, the influence of prompting is still observed: it influences verbosity, latency, result-only compliance, and sometimes even Pass@1. However, the severity of this influence depends on the task type, model, and evaluator. A neutral or polite tone might not always work, and the same goes for expert roles and formats. Thus, the influence of prompting has to be determined in specific setups, including the model, task type, evaluator, parser, and latency constraints.

\section{Purpose and Goal\label{sec:purpose_and_goal}}

This thesis is guided by four central questions regarding the size of prompt-framing effects, their impact on specific tasks and models, changes in observable behavior, and how relevant behavioral factors contribute to explaining differences in correctness. The detailed research questions are outlined in Section~\ref{sec:research-questions}.

This thesis makes four primary contributions:
\begin{itemize}
\tightlist
\item \textbf{A systematic experiment}: Full-factorial prompt-framing experiment combining tone, role framing, and prompt format across 14,400 LLM runs.
\item \textbf{Evaluating beyond correctness alone}: A joint evaluation of both correctness and observable behaviors, tracking multiple outcomes such as response length, speed, and compliance alongside accuracy \cite{liang2023helm}.
\item \textbf{Model- and task-stratified analysis}: Reporting results separately for each model and task category, showing that prompt effects are smaller than model effects but still meaningful for output behavior and engineering reliability.
\item \textbf{Mediation and sensitivity analysis}: Examining when prompt effects are associated with verbosity, compliance, and runtime, and when pooled averages are dominated by a single strong model.
\end{itemize}

\section{Thesis Outline\label{sec:outline}}

This thesis is organized into seven chapters as follows:

\textbf{Chapter \ref{cha:chapter1}} (Introduction) presents the motivation, research objectives, thesis scope, and primary contributions.

\textbf{Chapter \ref{cha:chapter2}} (Theoretical Background and Related Work) surveys existing literature on prompt engineering, prompt sensitivity, LLM code generation, structured prompts, and reasoning evaluation.

\textbf{Chapter \ref{cha:chapter3}} (Experimental Design) details the research questions, factorial prompt matrix, evaluated local and cloud models, task datasets, and execution environment.

\textbf{Chapter \ref{cha:chapter4}} (Statistical and Analytical Strategy) describes the mixed-effects logistic regression models, linear probability models, multiple comparison corrections, and behavioral mediation framework.

\textbf{Chapter \ref{cha:chapter5}} (Results and Empirical Analysis) presents empirical findings across model groups, correctness metrics, token usage, latency, compliance, and failure mode audits.

\textbf{Chapter \ref{cha:chapter6}} (Discussion and Threats to Validity) synthesizes key insights, practical prompt-policy guidelines, and internal, construct, statistical conclusion, and external validity threats.

\textbf{Chapter \ref{cha:chapter7}} (Conclusion and Outlook) summarizes the main findings, replication package, and outlines directions for future research.
